% \iffalse meta-comment
%
% letteredthm.dtx
% Copyright 2026 Gorkem Aksaray <aksarayg@tcd.ie>
%
% This work may be distributed and/or modified under the
% conditions of the LaTeX Project Public License, either version 1.3c
% of this license or (at your option) any later version.
% The latest version of this license is in
%   https://www.latex-project.org/lppl.txt
% and version 1.3c or later is part of all distributions of LaTeX
% version 2008 or later.
%
% This work has the LPPL maintenance status `maintained'.
%
% The Current Maintainer of this work is Gorkem Aksaray.
%
% This work consists of the files letteredthm.dtx, letteredthm.ins,
% and the derived files letteredthm.sty and letteredthm.pdf.
%
% \fi
%
% \iffalse
%<*driver>
\ProvidesFile{letteredthm.dtx}
%</driver>
%<package>\NeedsTeXFormat{LaTeX2e}[2021-11-15]
%<package>\ProvidesExplPackage{letteredthm}{2026-11-23}{1.0.0}
%<package>  {Lettered sub-numbering for theorem-like environments}
%<*driver>
    [2026/11/23 v1.0.0 Lettered sub-numbering for theorem-like environments]
%</driver>
%
%<*driver>
\documentclass{ltxdoc}
\usepackage[T1]{fontenc}
\usepackage{lmodern}
\usepackage[indent=15pt]{parskip}
\usepackage[dvipsnames]{xcolor}
\usepackage{microtype}
\usepackage[hyperindex=false]{hyperref}
\hypersetup{
  colorlinks,
  linkcolor=Red,
  urlcolor=RoyalBlue,
  citecolor=RoyalBlue,
  pdftitle={The letteredthm package},
  pdfauthor={Gorkem Aksaray},
  pdfsubject={Lettered sub-numbering for theorem-like environments}}
\newcommand{\email}[1]{\href{mailto:#1}{\nolinkurl{#1}}}
\RecordChanges
% expl3 names contain _ and :, so treat them as letters in macro names
\makeatletter
\renewcommand*{\MakePrivateLetters}{\makeatletter\catcode`\_=11 \catcode`\:=11 }
\makeatother
\begin{document}
  \DocInput{letteredthm.dtx}
  \PrintChanges
\end{document}
%</driver>
% \fi
%
% \CheckSum{100}
%
% \CharacterTable
%  {Upper-case    \A\B\C\D\E\F\G\H\I\J\K\L\M\N\O\P\Q\R\S\T\U\V\W\X\Y\Z
%   Lower-case    \a\b\c\d\e\f\g\h\i\j\k\l\m\n\o\p\q\r\s\t\u\v\w\x\y\z
%   Digits        \0\1\2\3\4\5\6\7\8\9
%   Exclamation   \!     Double quote  \"     Hash (number) \#
%   Dollar        \$     Percent       \%     Ampersand     \&
%   Acute accent  \'     Left paren    \(     Right paren   \)
%   Asterisk      \*     Plus          \+     Comma         \,
%   Minus         \-     Point         \.     Solidus       \/
%   Colon         \:     Semicolon     \;     Less than     \<
%   Equals        \=     Greater than  \>     Question mark \?
%   Commercial at \@     Left bracket  \[     Backslash     \\
%   Right bracket \]     Circumflex    \^     Underscore    \_
%   Grave accent  \`     Left brace    \{     Vertical bar  \|
%   Right brace   \}     Tilde         \~}
%
% \changes{v1.0.0}{2026/11/23}{Initial version}
%
% \GetFileInfo{letteredthm.dtx}
%
% \title{The \textsf{letteredthm} package}
% \author{Gorkem Aksaray \\ \email{aksarayg@tcd.ie}}
% \date{\fileversion~from \filedate}
%
% \maketitle
%
% \noindent Source repo:
% \url{https://github.com/gaksaray/letteredthm}.
%
% \section{Introduction}
%
% This is a \LaTeX\ package for theorem-like environments, such as
% hypotheses or propositions, that are numbered 1, 2, 3, \dots\ but can also
% be lettered in groups: 1, 2a, 2b, 3, 4a, 4b, 4c, 5, and so on. Any text,
% sections or floats can sit between the lettered statements of a group, so
% an argument can lead to Hypothesis~1a, continue, and then lead to
% Hypothesis~1b. With \textsf{cleveref}, references to plain and lettered
% statements are sorted into one list. The package needs nothing beyond the
% \LaTeX\ kernel and works with \textsf{amsthm} or with the kernel's own
% |\newtheorem|.
%
% \section{Usage}
%
% Load the package anywhere in the preamble:
% \begin{verbatim}
%   \usepackage{letteredthm}
% \end{verbatim}\vspace{-\baselineskip}
% It may be loaded before or after \textsf{cleveref} and \textsf{hyperref}.
%
% \DescribeMacro{\NewLetteredTheorem}
% \hfill|\NewLetteredTheorem|\marg{env}\marg{name}\marg{Singular}\marg{Plural}\hfill\null
%
% Defines the theorem-like environment \meta{env}, headed \meta{Singular},
% together with the commands |\start|\meta{name} and |\stop|\meta{name}.
% \meta{Singular} and \meta{Plural} also give the names that \textsf{cleveref}
% uses in references. The environment takes the current theorem style, so set
% it first with |\theoremstyle| if you use \textsf{amsthm}:
% \begin{verbatim}
%   \theoremstyle{definition}
%   \NewLetteredTheorem{hypothesis}{subhypotheses}{Hypothesis}{Hypotheses}
% \end{verbatim}\vspace{-\baselineskip}
% This defines the environment |hypothesis| and the commands
% |\startsubhypotheses| and |\stopsubhypotheses|.
%
% \DescribeMacro{\start\meta{name}}
% \DescribeMacro{\stop\meta{name}}
% \hfill|\start|\meta{name}\oarg{label}\quad\dots\quad|\stop|\meta{name}\hfill\null
%
% Outside a start/stop pair, each \meta{env} takes the next number. Between
% them, every \meta{env} shares the next number and gets a letter. The
% commands add no grouping or formatting, so anything can come between the
% lettered statements:
% \begin{verbatim}
%   For these reasons, we expect $Y$ to increase in $X$. Hence:
%   \startsubhypotheses[hyp:1]
%   \begin{hypothesis}\label{hyp:1a}
%     $Y$ increases in $X$.
%   \end{hypothesis}
%   Of course, we could make the same argument for $Z$. Hence:
%   \begin{hypothesis}\label{hyp:1b}
%     $Y$ increases in $Z$.
%   \end{hypothesis}
%   \stopsubhypotheses
%   \begin{hypothesis}\label{hyp:2}
%     ...
%   \end{hypothesis}
% \end{verbatim}\vspace{-\baselineskip}
% This gives Hypotheses 1a, 1b and 2. A new |\start|\meta{name} also closes
% any open group, so two groups can follow each other directly. The optional
% \meta{label} refers to the group as a whole: above, |\ref{hyp:1}| prints
% ``1'' and links to Hypothesis~1a.
%
% With \textsf{cleveref}, |\cref{hyp:2,hyp:1b,hyp:1a}| prints ``hypotheses
% 1a, 1b and 2'': plain and lettered statements are sorted into one list.
% Reference names set with |\crefname| or |\Crefname| take precedence over
% those given to |\NewLetteredTheorem|.
%
% \section{Notes}
%
% A group takes its number when its first statement starts, so an empty
% group uses up no number. If a group with a \meta{label} closes before any
% statement, the package warns that the label is undefined. Behind each
% \meta{env}, the package defines a second theorem environment,
% |sub|\meta{env}, for the lettered statements; it is not meant to be used
% directly. The package requires a \LaTeX\ format from November 2021 or
% later.
%
% \StopEventually{}
%
% \section{Implementation}
%
% The package is written in \textsf{expl3}, part of the \LaTeX\ kernel. The
% prefix of its internal names is |lthm|.
%    \begin{macrocode}
\skip_new:N \l__lthm_skip
\clist_new:N \g__lthm_envs_clist
%    \end{macrocode}
%
% \begin{macro}{\NewLetteredTheorem}
% Two theorem environments are created: \meta{env} for plain statements and
% |sub|\meta{env}, whose counter is reset by that of \meta{env}, for lettered
% ones. The original \meta{env} is saved, and \meta{env} is then redefined to
% call either environment depending on whether a group is open. The choice
% is made at the level of the environment's own commands rather than by
% nesting one environment inside another, which would lose \textsf{amsthm}'s
% handling of the paragraph that follows a statement.
%
% Each \meta{env} has three variables: whether a group is open, whether the
% open group still has to take its number, and the group's label.
%    \begin{macrocode}
\NewDocumentCommand \NewLetteredTheorem { m m m m }
  {
    \newtheorem {#1} {#3}
    \newtheorem {sub#1} {#3} [#1]
    \cs_set:cpn { thesub#1 } { \use:c { the#1 } \alph { sub#1 } }
    \bool_new:c { g__lthm_#1_group_bool }
    \bool_new:c { g__lthm_#1_pending_bool }
    \tl_new:c { g__lthm_#1_label_tl }
    \cs_new_eq:cc { __lthm_#1_plain: } {#1}
    \cs_new_eq:cc { __lthm_#1_endplain: } { end#1 }
    \cs_set_protected:cpn {#1}
      {
        \bool_if:cTF { g__lthm_#1_group_bool }
          { \__lthm_open_group:n {#1} \use:c { sub#1 } }
          { \use:c { __lthm_#1_plain: } }
      }
    \cs_set_protected:cpn { end#1 }
      {
        \bool_if:cTF { g__lthm_#1_group_bool }
          { \use:c { endsub#1 } }
          { \use:c { __lthm_#1_endplain: } }
      }
%    \end{macrocode}
% |\start|\meta{name} and |\stop|\meta{name} only set these variables, so
% they produce no output and no grouping.
%    \begin{macrocode}
    \exp_args:Nc \NewDocumentCommand { start#2 } { O{} }
      {
        \__lthm_check_unused:n {#1}
        \bool_gset_true:c { g__lthm_#1_group_bool }
        \bool_gset_true:c { g__lthm_#1_pending_bool }
        \tl_gset:cn { g__lthm_#1_label_tl } {##1}
      }
    \exp_args:Nc \NewDocumentCommand { stop#2 } { }
      {
        \__lthm_check_unused:n {#1}
        \bool_gset_false:c { g__lthm_#1_group_bool }
        \bool_gset_false:c { g__lthm_#1_pending_bool }
      }
%    \end{macrocode}
% Finally, the parent of |sub|\meta{env} and the reference names are recorded
% for the \textsf{cleveref} support below.
%    \begin{macrocode}
    \tl_const:cn { c__lthm_parent_sub#1_tl } {#1}
    \clist_gput_right:Nn \g__lthm_envs_clist { {#1} {#3} {#4} }
  }
%    \end{macrocode}
% \end{macro}
%
% \begin{macro}{\__lthm_open_group:n}
% Called as each lettered statement starts. The first statement of a group
% steps the parent counter, which also resets the letter, and writes the
% group's label, if any. The link target and the label are invisible marks
% that would keep the space above them from merging with the space
% \textsf{amsthm} adds above the statement, leaving a wider gap. So that
% space is taken out and put back after them; after a heading, a |\nobreak|
% also keeps the page from breaking between the heading and the statement.
%    \begin{macrocode}
\cs_new_protected:Npn \__lthm_open_group:n #1
  {
    \bool_if:cT { g__lthm_#1_pending_bool }
      {
        \bool_gset_false:c { g__lthm_#1_pending_bool }
        \mode_if_horizontal:T { \unskip \unskip \par }
        \skip_set:Nn \l__lthm_skip { \tex_lastskip:D }
        \skip_if_eq:nnF { \l__lthm_skip } { \c_zero_skip }
          { \skip_vertical:n { - \l__lthm_skip } }
        \refstepcounter {#1}
        \tl_if_empty:cF { g__lthm_#1_label_tl }
          { \exp_args:Nv \label { g__lthm_#1_label_tl } }
        \legacy_if:nT { @nobreak } { \nobreak }
        \skip_if_eq:nnF { \l__lthm_skip } { \c_zero_skip }
          { \skip_vertical:N \l__lthm_skip }
      }
  }
%    \end{macrocode}
% \end{macro}
%
% \begin{macro}{\__lthm_check_unused:n}
% Warns when a group closes while its label is still waiting for a statement.
%    \begin{macrocode}
\cs_new_protected:Npn \__lthm_check_unused:n #1
  {
    \bool_lazy_and:nnT
      { \bool_if_p:c { g__lthm_#1_pending_bool } }
      { ! \tl_if_empty_p:c { g__lthm_#1_label_tl } }
      {
        \exp_args:Nnnv \msg_warning:nnn { letteredthm } { unused-label }
          { g__lthm_#1_label_tl }
      }
  }
\msg_new:nnn { letteredthm } { unused-label }
  { Group~label~`#1'~is~undefined:~its~group~has~no~statements. }
%    \end{macrocode}
% \end{macro}
%
% \begin{macro}{\__lthm_cref_setup:}
% The \textsf{cleveref} support is set up at |\begin{document}|, so that the
% package can be loaded before or after \textsf{cleveref}. \textsf{cleveref}
% sorts references by their counters' values within their parent counters,
% but it finds parents only from a fixed list (section, figure, and so on).
% Without help it would sort 2a, 3a, 2b by their letters alone, so its
% |\cref@resetby| is extended to name the parent of each |sub|\meta{env}.
%    \begin{macrocode}
\AddToHook { begindocument / before }
  { \IfPackageLoadedTF { cleveref } { \__lthm_cref_setup: } { } }
\cs_new_protected:Npn \__lthm_cref_setup:
  {
    \clist_map_function:NN \g__lthm_envs_clist \__lthm_cref_names:n
    \cs_set_eq:Nc \__lthm_cref_resetby:nN { cref@resetby }
    \cs_set:cpn { cref@resetby } ##1##2
      {
        \__lthm_cref_resetby:nN {##1} ##2
        \tl_if_exist:cT { c__lthm_parent_##1_tl }
          { \tl_set_eq:Nc ##2 { c__lthm_parent_##1_tl } }
      }
  }
%    \end{macrocode}
% \end{macro}
%
% \begin{macro}{\__lthm_cref_names:n}
% \begin{macro}{\__lthm_cref_names:nnn}
% Sets the reference names of each \meta{env}, unless the user has already
% set them, and makes references to |sub|\meta{env} count as references to
% \meta{env}, so that both are listed together.
%    \begin{macrocode}
\cs_new_protected:Npn \__lthm_cref_names:n #1
  { \__lthm_cref_names:nnn #1 }
\cs_new_protected:Npn \__lthm_cref_names:nnn #1#2#3
  {
    \cs_if_exist:cF { Cref@#1@name } { \Crefname {#1} {#2} {#3} }
    \cs_if_exist:cF { cref@#1@name }
      { \crefname {#1} { \MakeLowercase {#2} } { \MakeLowercase {#3} } }
    \crefalias { sub#1 } {#1}
  }
%    \end{macrocode}
% \end{macro}
% \end{macro}
%
% \Finale
\endinput
